problem:  
Let \( X \) be an irreducible bounded symmetric domain of complex dimension \( n \), realized in its Harish–Chandra embedding with Bergman metric \( g_B \). Suppose \( \Omega \subset \mathbb{C}^n \) is a bounded, smoothly bounded, strictly pseudoconvex domain whose Bergman metric is \( g_\Omega \). Assume there exists a holomorphic isometric embedding  
\[
F : (\Omega, g_\Omega) \hookrightarrow (X, g_B).
\]  
Must \( \Omega \) necessarily be biholomorphic to a bounded symmetric domain (and thus itself Hermitian symmetric)?  
Equivalently, does the existence of a holomorphic Bergman isometric embedding into a Hermitian symmetric space force the source domain to possess global symmetric structure?

Why it is a "good" problem:  
This problem proposes a new rigidity principle linking the intrinsic complex geometry of strictly pseudoconvex domains to the global symmetry of Hermitian symmetric spaces. It advances classical rigidity theory (Calabi, Siu, Mok), which typically assumes the source domain is already symmetric (e.g., the unit ball), by probing whether symmetry can be *induced* purely from the isometric property of Bergman metrics. The question highlights a sharp geometric contrast: strictly pseudoconvex domains exhibit holomorphic sectional curvature strongly negative near the boundary, while bounded symmetric domains—especially those of higher rank—carry richer curvature anisotropy reflecting their Lie-theoretic structure. Determining whether an isometric embedding can reconcile these curvature models would require deep interaction among several complex variables, Kähler geometry, CR boundary theory, and representation theory.  

The problem is elegantly simple to state yet conceals significant analytic and geometric depth. It bridges local analytic data (the Bergman metric curvature of a strictly pseudoconvex domain) and global algebraic structure (the symmetric property of \( X \)), and its resolution would clarify how completely the Bergman metric encodes global symmetry. Thus, it fully embodies the “narrow entrance, profound depth” and cross‑disciplinary power that characterize a good research‑level mathematical problem.